\SGAsection{附录 I}{Verdier 的一个反例}
\footnote{除少数行文细节外，本附录重现了 J.-L. Verdier 致编者的两封信（1967 年 9 月 30 日及 10 月 11 日），以回答编者提出的一个问题。}
\SourcePage{205}
\SGAsection{0}{引言}
若 $S$ 是一个\Term{scheme}{概形}，以 $\operatorname{Mod}(S)$（相应地 $\operatorname{Qcoh}(S)$）记 $\mathcal O_S$-\Term{module}{模}（相应地\Term{quasi-coherent}{拟凝聚} $\mathcal O_S$-\Term{module}{模}）的\Term{category}{范畴}。本附录旨在提醒读者注意\Term{category}{范畴} $D(\operatorname{Qcoh}(S))$ 的某些病态性质，特别是证明下列事实。

\textbf{0.1.} 设 $X$ 是环为 $A$ 的\Term{affine-scheme}{仿射概形}，$I$ 是一个\Term{injective-module}{内射} $A$-\Term{module}{模}。已知（[H] II 7.14），若 $A$ 是\Term{noetherian}{诺特} 的，则\Term{sheaf}{层} $\widetilde I$ 是一个\Term{injective-module}{内射} $\mathcal O_X$-\Term{module}{模}。若 $A$ 不是\Term{noetherian}{诺特} 的，则 $\widetilde I$ 未必\Term{injective-module}{内射}，甚至未必\Term{flasque}{松弛}，即使 $X$ 的底\Term{space}{空间}是\Term{noetherian}{诺特} 的且 \Term{krull-dimension}{Krull 维数}\Term{finite}{有限}。此外，$\widetilde I$ 在 $X$ 的\Term{open-subset}{开集} $U$ 上的\Term{restriction}{限制}\Term{generic}{一般}也不是 $\operatorname{Qcoh}(U)$ 的\Term{injective-object}{内射对象}。

\textbf{0.2.} 设 $f:X\to Y$ 是\Term{quasi-compact}{拟紧}且\Term{quasi-separated}{拟分离}的\Term{scheme}{概形}之间的\Term{morphism}{态射}。\Term{functor}{函子} $\mathrm Rf_*:D^+(X)\to D^+(Y)$ 在 $D^+(\operatorname{Qcoh}(X))$（由第 II 讲 3.5，这是 $D^+(X)$ 的\Term{full-subcategory}{全子范畴}）上诱导的\Term{functor}{函子}，\Term{generic}{一般}不是 $f_*:\operatorname{Qcoh}(X)\to\operatorname{Mod}(Y)$ 的\Term{right-derived-functor}{右导出函子}。更确切地说，$\operatorname{Qcoh}(X)$ 的\Term{injective-object}{内射对象}对于 $f_*:\operatorname{Mod}(X)\to\operatorname{Mod}(Y)$ 未必是\Term{acyclic}{无上同调}的。

\textbf{0.3.} 设 $S$ 是\Term{quasi-compact}{拟紧}且\Term{quasi-separated}{拟分离}的\Term{scheme}{概形}。$\operatorname{Qcoh}(S)$ 的一个\Term{injective-object}{内射对象}对于\Term{coherator}{凝聚化函子}（第 II 讲 3.2）未必\Term{acyclic}{无上同调}。\Term{functor}{函子} $\delta:D^b(\operatorname{Qcoh}(S))\to D^b(S)$（第 II 讲 3.0）\Term{generic}{一般}不是\Term{fully-faithful}{全忠实}的。

\SGAsection{1}{反例}
先回顾下列\Term{lemma}{引理}（SGA 2 II 9）。
\begin{SGAresult}{\Term{lemma}{引理}}{1.1}
设 $A$ 是一个环，$X=\operatorname{Spec}(A)$，$\underline f=(f_\alpha)$ 是 $A$ 中的一个\Term{finite}{有限}元素族，$Y$ 是由 $\underline f(x)=0$ 定义的\Term{closed-subset}{闭集}。对 $i>0$，下列条件\Term{equivalence}{等价}：
\begin{enumerate}
\item[(i)] 对每个\Term{injective-module}{内射} $A$-\Term{module}{模} $I$，有 $H^i_Y(X,\widetilde I)=0$；
\item[(ii)] \Term{pro-object}{pro-对象} $n\mapsto H_i(\underline f^{\,n},A)$ 为零，即对每个 $n$ 都有 $n'\geq n$，使得
\SourcePage{206}
\[H_i(\underline f^{\,n'},A)\longrightarrow H_i(\underline f^{\,n},A)\]
为零。
\end{enumerate}
\end{SGAresult}
\textbf{1.2.} 我们用 1.1 证明 0.1。设 $k$ 是一个域，$B=k[[X,Y]]$，$\underline f$ 是元素族 $(x,y)$，$M=\bigoplus_{n\geq0}B/(\underline f^{\,n})$，$A$ 是 $B$-\Term{algebra}{代数} $D_B(M)$（EGA $0_{\mathrm{IV}}$ 18.2.3）。显然有
\[H_2^A(\underline f^{\,n},A)\simeq H_2^B(\underline f^{\,n},A).\footnote{上标 $A$ 或 $B$ 指明计算 Koszul 复形所相对于的环。}\]
另一方面有
\[H_2^B(\underline f^{\,n},A)\simeq H_2^B(\underline f^{\,n},B)\oplus H_2^B(\underline f^{\,n},M),\]
但 $H_2^B(\underline f^{\,n},B)=0$，因为 $\underline f$ 是\Term{regular-sequence}{正则序列}（EGA III 1.1.4 及 $0_{\mathrm{IV}}$ 15.1.20），所以
\[H_2^B(\underline f^{\,n},A)\simeq H_2^B(\underline f^{\,n},M).\]
最后（EGA III 1.1.3.3）有
\[H_2^B(\underline f^{\,n},M)\simeq H_B^0(\underline f^{\,n},M),\]
其中 $H_B^0(\underline f^{\,n},M)$ \Term{canonical}{典范}地等同于 $M$ 中被 $(\underline f^{\,n})$ 消去的\Term{submodule}{子模} ${}_{(\underline f^{\,n})}M$。因此最终得到
\[H_2^A(\underline f^{\,n},A)\simeq{}_{(\underline f^{\,n})}M.\tag{*}\]
由于过渡\Term{morphism}{态射} $H_2^A(\underline f^{\,n'},A)\to H_2^A(\underline f^{\,n},A)$ 由乘以 $\underline f^{\,n'-n}$ 诱导（EGA III 1.1.6），利用 (*) 立即可见\Term{pro-object}{pro-对象} $n\mapsto H_2^A(\underline f^{\,n},A)$ 不为零。由 1.1 可知存在一个\Term{injective-module}{内射} $A$-\Term{module}{模} $I$，使 $H_Y^2(X,\widetilde I)\ne0$，从而 $\widetilde I$ 不是\Term{flasque}{松弛层}（SGA 2 I 2.12）。顺便指出，\Term{projection}{投影} $\operatorname{Spec}(A)\to\operatorname{Spec}(B)$ 是一个\Term{homeomorphism}{同胚}，
\SourcePage{207}
所以 $X$ 的底\Term{space}{空间}是\Term{dimension}{维数}为 $2$ 的\Term{noetherian}{诺特} \Term{space}{空间}。取一个\Term{injective-module}{内射} $A$-\Term{module}{模} $I$，使 $H_Y^2(X,\widetilde I)\ne0$，并以 $U$ 记 $Y$ 在 $X$ 中的开补集。由 SGA 2 II 4，有
\[H_Y^2(X,\widetilde I)\simeq H^1(U,\widetilde I)\ne0.\]
而由于 $U$ 是\Term{separated}{分离}的，\Term{functor}{函子} $D^+(\operatorname{Qcoh}(U))\to D^+(U)$ 是\Term{fully-faithful}{全忠实}的（第 II 讲 3.5），所以 $\widetilde I|U$ 在 $\operatorname{Qcoh}(U)$ 中不是\Term{injective-object}{内射对象}。这就证明了断言 0.1。

以 $Z$ 记沿 $U$ \Term{gluing}{粘合}两个 $X$ 的副本所得到的\Term{quasi-compact}{拟紧}且\Term{quasi-separated}{拟分离}\Term{scheme}{概形}，以 $j_1,j_2$ 记 $X$ 到 $Z$ 的两个\Term{canonical}{典范}\Term{open-immersion}{开浸入}，于是有\Term{cartesian-square}{笛卡尔方块}
\[\begin{tikzcd}
U\arrow[r,"i"]\arrow[d,"i"']&X\arrow[d,"j_1"]\\
X\arrow[r,"j_2"']&Z.
\end{tikzcd}\]
有
\[j_2^*\mathrm R^1j_{1*}(\widetilde I)\simeq\mathrm R^1i_*(i^*\widetilde I),\]
故 $H^0(X,\mathrm R^1j_{1*}(\widetilde I))\simeq H^1(U,\widetilde I)\ne0$，从而 $\mathrm R^1j_{1*}(\widetilde I)\ne0$，这就证明了 0.2。

另一方面，由第 II 讲 3.4 及 3.5，有
\[\mathrm RQ\,\mathrm Rj_{1*}(\widetilde I)\simeq j_{1*}(\widetilde I).\]
由此得到一个\Term{spectral-sequence}{谱序列}，其低次项\Term{exact-sequence}{正合列}给出\Term{isomorphism}{同构}
\[\mathrm R^1j_{1*}(\widetilde I)\simeq\mathrm R^2Q(j_{1*}\widetilde I).\]
\Term{sheaf}{层} $j_{1*}\widetilde I$ 在 $\operatorname{Qcoh}(Z)$ 中是\Term{injective-object}{内射对象}，所以这就证明了 0.3 的第一部分。
\SourcePage{208}
最后，设 $E\in\operatorname{ob}\operatorname{Qcoh}(Z)$。有\Term{duality}{对偶}\Term{isomorphism}{同构}
\[\mathrm R\operatorname{Hom}_{\mathrm{Mod}}(E,\mathrm Rj_{1*}(\widetilde I))\simeq\mathrm R\operatorname{Hom}_{\mathrm{Mod}}(j_1^*E,\widetilde I),\]
但由于 $X$ 是\Term{affine-scheme}{仿射概形}，由第 II 讲 3.5 有
\[\mathrm R\operatorname{Hom}_{\mathrm{Mod}}(j_1^*E,\widetilde I)\simeq\mathrm R\operatorname{Hom}_{\mathrm{Qcoh}}(j_1^*E,\widetilde I)\simeq\operatorname{Hom}(j_1^*E,\widetilde I).\]
所以
\[\mathrm R\operatorname{Hom}_{\mathrm{Mod}}(E,\mathrm Rj_{1*}(\widetilde I))\simeq\operatorname{Hom}(j_1^*E,\widetilde I),\]
由此得到一个\Term{spectral-sequence}{谱序列}，其低次项\Term{exact-sequence}{正合列}给出\Term{isomorphism}{同构}
\[\operatorname{Ext}_{\mathrm{Mod}}^2(E,j_{1*}\widetilde I)\simeq\operatorname{Hom}(E,\mathrm R^1j_{1*}\widetilde I).\]
因此，$E=\mathrm R^1j_{1*}(\widetilde I)$ 的\Term{identity-morphism}{恒等态射}给出 $\operatorname{Ext}_{\mathrm{Mod}}^2(E,j_{1*}\widetilde I)$ 中的一个非零元素；所以\Term{functor}{函子} $D^b(\operatorname{Qcoh}(Z))\to D^b(Z)$ 不是\Term{fully-faithful}{全忠实}的。这就完成了 0.3 的证明。

\SourcePage{209}
\SGAsection{附录 II}{相对椭圆复形的解析指标的定义}
本附录将说明，\Term{perfect-complex}{完美复形}的概念能对\Term{continuous}{连续}的\Term{elliptic-operator}{椭圆微分算子}族，或更\Term{generic}{一般}的“\Term{relative-elliptic-complex}{相对椭圆复形}”的“\Term{analytic-index}{解析指标}”给出简单而自然的定义。对于熟悉这一领域的读者，指出这里的定义较经典定义（例如见 [2]、[17]）的优点在于，它不需要对所考虑的\Term{elliptic-complex}{椭圆复形}作\Term{deformation}{形变}。

第 1、2 节回顾经典的“\Term{mixed-structure}{混合结构}”概念\footnote{建议熟悉这些内容的读者跳过这两节。}。当然，我们只给出在第 3 节陈述并证明\Term{center}{中心}\Term{theorem}{定理}所需的最少定义和结果。该\Term{theorem}{定理}更确切地说是：\Term{relative-elliptic-complex}{相对椭圆复形}沿到\Term{base}{基底}的\Term{projection}{投影}（假定\Term{proper}{固有}，\Term{base}{基底}假定\Term{locally-compact}{局部紧}）的\Term{direct-image}{直像}是\Term{perfect-complex}{完美复形}。当\Term{base}{基底} $S$ 紧时，由第 II 讲 2.3.2 的整体化论述即可把\Term{analytic-index}{解析指标}定义为 $K(S)$（$S$ 上\Term{finite-rank}{有限秩}复\Term{vector-bundle}{向量丛}的“Grothendieck 群”）的元素。

以下始终约定：若 $S$ 是\Term{topological-space}{拓扑空间}，以 $\mathcal O_S$ 记 $S$ 上实值\Term{continuous}{连续}函数的\Term{sheaf}{层}。
\SGAsection{1}{混合结构的基本论述}
本节固定一个\Term{topological-space}{拓扑空间} $S$。
\textbf{1.1. \Term{mixed-manifold}{混合流形}的定义。}
设 $V$ 是 $\mathbb R^n$ 中的\Term{open-subset}{开集}。对 $r\in\mathbb N$，以 $C^r(V)$ 记 $V$ 上实值 $C^r$ 函数的 Fréchet \Term{space}{空间}。置 $C^\infty(V)=\varprojlim_r C^r(V)$。
设 $U$ 是 $S$ 中的\Term{open-subset}{开集}。对 $r\in\mathbb N\cup\{\infty\}$，以 ${}^r\mathcal O_{U\times V}$ 记 $U\times V$ 上由下式定义的\Term{sheaf-of-rings}{环层}：
\SourcePage{210}
\[\Gamma(U'\times V',{}^r\mathcal O_{U\times V})=C(U',C^r(V')),\]
其中 $U'$（相应地 $V'$）是 $U$（相应地 $V$）中的\Term{open-subset}{开集}，$C(U',C^r(V'))$ 表示从 $U'$ 到 $C^r(V')$ 的\Term{continuous-map}{连续映射}\Term{space}{空间}。有\Term{canonical}{典范}\Term{injection}{单射}
\[\operatorname{pr}_1^{-1}(\mathcal O_U)\longrightarrow{}^r\mathcal O_{U\times V},\]
它使 $(U\times V,{}^r\mathcal O_{U\times V})$ 成为 $U$ 上的\Term{ringed-space}{赋环空间}。

称 $S$ 上的\Term{ringed-space}{赋环空间} $X$ 为一个 $C^r$ \Term{mixed-manifold}{混合流形}（$r\in\mathbb N\cup\{\infty\}$），若 $X$ 作为 $S$ 上的\Term{ringed-space}{赋环空间}\Term{local-scheme}{局部}\Term{isomorphic}{同构}于 $(U\times V,{}^r\mathcal O_{U\times V})$ 型的\Term{space}{空间}。（若 $S$ 只有一个点，$S$ 上的 $C^r$ \Term{mixed-manifold}{混合流形}就是通常意义下的 $C^r$ \Term{manifold}{流形}。）

设 $f:X\to S$ 是一个 $C^r$ \Term{mixed-manifold}{混合流形}。对 $s\in S$，\Term{fiber}{纤维} $X_s=f^{-1}(s)$ 赋以\Term{sheaf}{层} $i_s^{-1}(\mathcal O_X)/I_s$ 后是一个 $C^r$ \Term{manifold}{流形}，其中 $i_s:X_s\to X$ 是\Term{inclusion}{包含映射}，$I_s$ 是在 $X_s$ 上“为零”的\Term{section}{截面}组成的\Term{ideal}{理想}。

\textbf{1.2. “相对”\Term{differential-operator}{微分算子}。}
仿照\Term{algebraic-geometry}{代数几何}（EGA IV 16，SGA 3 VII），容易对\Term{mixed-manifold}{混合流形}进行\Term{infinitesimal-study}{无穷小研究}。这里只简要回顾所需的概念。

\textbf{1.2.1.} 设 $f:X\to S$ 是一个 $C^\infty$ \Term{mixed-manifold}{混合流形}（1.1）。对每个 $n\in\mathbb N$，定义 $X$ 上的 $n$ 阶\Term{principal-parts}{主部}\Term{sheaf}{层} $\mathcal P^n_{X/S}$。它通过 $X\times_S X$ 到 $X$ 的两个\Term{projection}{投影} $\operatorname{pr}_1,\operatorname{pr}_2$，以两种不同方式成为一个\Term{augmentation}{增广}到 $\mathcal O_X$ 的 $\mathcal O_X$-\Term{algebra}{代数}，作为 $\mathcal O_X$-\Term{module}{模}是\Term{locally-free-finite-type}{有限型局部自由}的。未说明所考虑的 $\mathcal P^n_{X/S}$ 的 $\mathcal O_X$-\Term{module}{模}结构时，均指来自 $\operatorname{pr}_1$ 的结构。有\Term{canonical}{典范}\Term{isomorphism}{同构}
\SourcePage{211}
\[\mathcal P^0_{X/S}\simeq\mathcal O_X,\qquad \mathcal P^1_{X/S}\simeq\mathcal O_X\oplus\Omega^1_{X/S},\]
其中 $\Omega^1_{X/S}$ 是 $X$ 相对于 $S$ 的\Term{relative-differentials}{相对微分}\Term{sheaf}{层}，也称\Term{relative-cotangent-sheaf}{相对余切层}。以 $T_{X/S}$ 记 $\Omega^1_{X/S}$ 的\Term{dual}{对偶}，称为\Term{relative-tangent-sheaf}{相对切层}。\Term{sheaf}{层} $T_{X/S}$（相应地 $\Omega^1_{X/S}$）是 $X$ 上一个\Term{finite-rank}{有限秩}\Term{vector-bundle}{向量丛}的\Term{section}{截面}\Term{sheaf}{层}；这个丛称为\Term{relative-tangent-bundle}{相对切丛}（相应地\Term{relative-cotangent-bundle}{相对余切丛}），它是 $X$ 上的一个 $S$-\Term{mixed-manifold}{混合流形}。

将 $\mathcal O_X$ 的\Term{section}{截面}送到其 $n$ 阶\Term{principal-parts}{主部}的 $f^{-1}(\mathcal O_S)$-\Term{linear}{线性}\Term{homomorphism}{同态}（对应于 $\operatorname{pr}_2$ 定义的 $\mathcal O_X$-\Term{algebraic-structure}{代数结构}）
\[d^n_{X/S}:\mathcal O_X\longrightarrow\mathcal P^n_{X/S}\]
称为 $X$ 相对于 $S$ 的 $n$ 阶\Term{universal-differential-operator}{泛微分算子}。

\textbf{1.2.2.} 更\Term{generic}{一般}地，若 $E$ 是\Term{locally-free-finite-type}{有限型局部自由} $\mathcal O_X$-\Term{module}{模}，可考虑 $E$ 的 $n$ 阶\Term{principal-parts}{主部}\Term{sheaf}{层} $\mathcal P^n_{X/S}(E)=\mathcal P^n_{X/S}\otimes_{\mathcal O_X}E$（\Term{tensor-product}{张量积}用 $\mathcal P^n_{X/S}$ 上的第二个 $\mathcal O_X$-\Term{module}{模}结构定义），并有一个\Term{universal-differential-operator}{泛微分算子}
\[d^n_{X/S}(E):E\longrightarrow\mathcal P^n_{X/S}(E).\]
设 $E,F$ 是\Term{locally-free-finite-type}{有限型局部自由} $\mathcal O_X$-\Term{module}{模}。按定义，从 $E$ 到 $F$ 的阶数 $\leq n$ 的\Term{differential-operator}{微分算子}是一个 $f^{-1}(\mathcal O_X)$-\Term{linear}{线性}\Term{homomorphism}{同态} $E\to F$，它经由 $d^n_{X/S}(E)$ \Term{resolution}{分解}为一个 $\mathcal O_X$-\Term{linear}{线性}\Term{homomorphism}{同态} $\mathcal P^n_{X/S}(E)\to F$（此时\Term{resolution}{分解}唯一）。换言之，\Term{operator}{算子} $d^n_{X/S}(E)$ 给出\Term{bijection}{双射}
\[\operatorname{Hom}_{\mathcal O_X}(\mathcal P^n_{X/S}(E),F)\xrightarrow{\sim}\operatorname{Diff}^n(E,F),\tag{1.2.2.1}\]
其中 $\operatorname{Diff}^n(E,F)$ 表示从 $E$ 到 $F$ 的阶数 $\leq n$ 的\Term{differential-operator}{微分算子}\Term{set}{集合}。
\SourcePage{212}
\textbf{1.2.3.} 我们还要考虑具有复系数的\Term{differential-operator}{微分算子}。以 $\mathcal O_S\otimes\mathbb C$（相应地 $\mathcal O_X\otimes\mathbb C$）记 $\mathcal O_S$（相应地 $\mathcal O_X$）的复化。设 $E,F$ 是\Term{locally-free-finite-type}{有限型局部自由} $(\mathcal O_X\otimes\mathbb C)$-\Term{module}{模}。所谓从 $E$ 到 $F$ 的阶数 $\leq n$ 的\Term{differential-operator}{微分算子}，是指一个 $f^{-1}(\mathcal O_S\otimes\mathbb C)$-\Term{linear}{线性}\Term{homomorphism}{同态} $E\to F$，它经由 $d^n_{X/S}(E)$ \Term{resolution}{分解}为 $(\mathcal O_X\otimes\mathbb C)$-\Term{linear}{线性}\Term{homomorphism}{同态} $\mathcal P^n_{X/S}(E)\to F$。若以 $\operatorname{Diff}_{\mathbb C}^n(E,F)$ 记这些\Term{operator}{算子}的\Term{set}{集合}，则\Term{operator}{算子} $d^n_{X/S}(E)$ 定义\Term{bijection}{双射}
\[\operatorname{Hom}_{\mathcal O_X\otimes\mathbb C}(\mathcal P^n_{X/S}(E),F)\xrightarrow{\sim}\operatorname{Diff}_{\mathbb C}^n(E,F).\tag{1.2.3.1}\]
\textbf{1.2.4.} 对 $n\geq1$，有\Term{canonical}{典范}\Term{exact-sequence}{正合列}
\[0\longrightarrow S^n(\Omega^1_{X/S})\longrightarrow\mathcal P^n_{X/S}\longrightarrow\mathcal P^{n-1}_{X/S}\longrightarrow0.\footnote{$S^n(\ )$ 表示第 $n$ 次\Term{symmetric-power}{对称幂}。}\]
因此与一个\Term{locally-free-finite-type}{有限型局部自由} $\mathcal O_X$-\Term{module}{模}（相应地 $(\mathcal O_X\otimes\mathbb C)$-\Term{module}{模}）$E$ 在 $\mathcal O_X$ 上取\Term{tensor-product}{张量积}\footnote{对于 $\mathcal P^n_{X/S}$ 及 $\mathcal P^{n-1}_{X/S}$ 的第二个 $\mathcal O_X$-\Term{module}{模}结构。}，得到 $\mathcal O_X$-\Term{module}{模}（相应地 $(\mathcal O_X\otimes\mathbb C)$-\Term{module}{模}）的\Term{exact-sequence}{正合列}
\[0\longrightarrow S^n(\Omega^1_{X/S})\otimes_{\mathcal O_X}E\longrightarrow\mathcal P^n_{X/S}(E)\longrightarrow\mathcal P^{n-1}_{X/S}(E)\longrightarrow0.\tag{1.2.4.1}\]
设 $F$ 是\Term{locally-free-finite-type}{有限型局部自由} $\mathcal O_X$-\Term{module}{模}（相应地 $(\mathcal O_X\otimes\mathbb C)$-\Term{module}{模}）。对 1.2.4.1 应用\Term{functor}{函子} $\operatorname{Hom}_{\mathcal O_X}(\ ,F)$（相应地 $\operatorname{Hom}_{\mathcal O_X\otimes\mathbb C}(\ ,F)$），并利用 1.2.2.1 及 1.2.4.1 中的等同，得到\Term{exact-sequence}{正合列}
\[0\to\operatorname{Diff}^{n-1}(E,F)\to\operatorname{Diff}^{n}(E,F)\xrightarrow{\sigma_n}\operatorname{Hom}_{\mathcal O_X}(S^n(\Omega^1_{X/S})\otimes_{\mathcal O_X}E,F),\tag{1.2.4.2}\]
相应地为
\[0\to\operatorname{Diff}_{\mathbb C}^{n-1}(E,F)\to\operatorname{Diff}_{\mathbb C}^{n}(E,F)\to\operatorname{Hom}_{\mathcal O_X\otimes\mathbb C}(S^n(\Omega^1_{X/S})\otimes_{\mathcal O_X}E,F).\]
\SourcePage{213}
（若 $X$ \Term{paracompact}{仿紧}且 $\mathcal O_X$ 是\Term{soft-sheaf}{软层}，见第 II 讲 2.3.1，则可在 1.2.4.2 右端补上一个零。）\Term{homomorphism}{同态} $\sigma_n$ 称为\Term{symbol}{符号}；它把一个从 $E$ 到 $F$ 的阶数 $\leq n$ 的\Term{differential-operator}{微分算子}送到一个“从 $\Omega^1_{X/S}$ 到 $\operatorname{Hom}_{\mathcal O_X}(E,F)$（相应地 $\operatorname{Hom}_{\mathcal O_X\otimes\mathbb C}(E,F)$）的 $n$ 次\Term{homogeneous-polynomial-map}{齐次多项式映射}”。也可将它看作一个\Term{homomorphism}{同态} $p^*(E)\to p^*(F)$，其中 $p$ 是\Term{relative-cotangent-bundle}{相对余切丛}到 $X$ 的\Term{projection}{投影}（1.2.1），并满足某个\Term{homogeneous}{齐次}性条件。

设 $E,F,G$ 是\Term{locally-free-finite-type}{有限型局部自由} $\mathcal O_X$-\Term{module}{模}，$u\in\operatorname{Diff}^m(E,F)$、$v\in\operatorname{Diff}^n(F,G)$，从而 $vu\in\operatorname{Diff}^{m+n}(E,G)$。容易验证关系
\[\sigma_{m+n}(vu)=\sigma_n(v)\sigma_m(u).\tag{1.2.4.3}\]
（复系数\Term{differential-operator}{微分算子}也有相同关系。）特别地，若两个\Term{differential-operator}{微分算子}的复合为零，它们的\Term{symbol}{符号}的复合也为零。

\textbf{1.2.5.} 设 $E$ 是 $X$ 相对于 $S$ 的复系数\Term{differential-operator}{微分算子}组成的\Term{bounded}{有界}\Term{complex}{复形}（1.2.3），其中 $d^i:E^i\to E^{i+1}$ 的阶数 $\leq n_i$。由 1.2.4.3，有
\[\sigma_{n_{i+1}}(d^{i+1})\sigma_{n_i}(d^i)=0.\]
称 $E$ 是阶数为 $\underline n=(n_i)$ 的\Term{elliptic-complex}{椭圆复形}，若\Term{complex}{复形} $(p^*(E^i),\sigma_{n_i}(d^i))$（$p$ 是\Term{cotangent-bundle}{余切丛}到 $X$ 的\Term{projection}{投影}）在\Term{zero-section}{零截面}之外是\Term{acyclic}{无上同调}的。

当 $S$ 只有一个点时，上述 $X$ 上的\Term{elliptic-complex}{椭圆复形}概念与通常的概念一致，后者在 [4] 或 [16] 中有详细研究。

\textbf{1.2.6.} 设 $s\in S$，$i_s$ 是\Term{fiber}{纤维} $X_s$ 到 $X$ 的\Term{inclusion}{包含映射}（看作\Term{ringed-space}{赋环空间}的\Term{morphism}{态射}，见 1.1）。对 $n\in\mathbb N$，有\Term{canonical}{典范}\Term{isomorphism}{同构}
\SourcePage{214}
\[i_s^*(\mathcal P^n_{X/S})\simeq\mathcal P^n_{X_s},\]
从而对\Term{locally-free-finite-type}{有限型局部自由} $\mathcal O_X$-\Term{module}{模} $E,F$ \Term{finite}{有限}制\Term{homomorphism}{同态}
\[i_s^*:\operatorname{Diff}^n(E,F)\longrightarrow\operatorname{Diff}^n(i_s^*(E),i_s^*(F))\]
（复系数\Term{operator}{算子}也有类似的\Term{homomorphism}{同态}）。设 $E$ 是 $X$ 上复系数\Term{differential-operator}{微分算子}组成的\Term{bounded}{有界}\Term{complex}{复形}；$E$ 是\Term{elliptic-complex}{椭圆复形}（1.2.5）的充要条件是：对每个 $s\in S$，$i_s^*(E)$ 是\Term{elliptic-complex}{椭圆复形}（留给读者验证）。

\SGAsection{2}{刚性引理}
\textbf{2.0.} “众所周知”，紧\Term{differentiable-manifold}{微分流形}是“\Term{rigid}{刚性}”的，即它们除了\Term{trivial-deformation}{平凡形变}外没有其他\Term{continuous}{连续}\Term{deformation}{形变}。不过，由于这个结果在第 3 节起关键\Term{action}{作用}，而可参阅的文献似乎相当少，给出准确陈述及若干证明或许并非多余。

固定一个\Term{topological-space}{拓扑空间} $S$。所讨论的\Term{mixed-manifold}{混合流形}均假定为 $C^r$ 类，$r\geq1$（也可 $r=\infty$）。

\textbf{2.1. 记号。} 设 $X,Y$ 是 $S$-\Term{mixed-manifold}{混合流形}。以 $\operatorname{Hom}(X,Y)$ 记 $X$ 到 $Y$ 的 $S$-\Term{morphism}{态射}\Term{set}{集合}，以 $\mathcal Hom(X,Y)$ 记 $X$ 上从 $X$ 到 $Y$ 的 $S$-\Term{morphism}{态射}的芽组成的\Term{sheaf}{层}（由 $\mathcal Hom(X,Y)(U)=\operatorname{Hom}(U,Y)$ 定义，$U$ 为 $X$ 的\Term{open-subset}{开集}）。
设 $s\in S$，以 $i_s$ 记 $X_s$ 到 $X$ 的\Term{inclusion}{包含映射}。有\Term{canonical}{典范}的诱导\Term{homomorphism}{同态}
\[i_s^{-1}\mathcal Hom(X,Y)\longrightarrow\mathcal Hom(X_s,Y_s),\tag{2.1.1}\]
\SourcePage{215}
它是\Term{epimorphism}{满射}，这一点立即可见。当 $Y=S\times\mathbb R$ 时，\Term{sheaf}{层} $\mathcal Hom(X,Y)$ \Term{canonical}{典范}地等同于 $\mathcal O_X$，而\Term{epimorphism}{满射} 2.1.1 就等同于 1.1 中考虑的\Term{epimorphism}{满射}
\[i_s^{-1}(\mathcal O_X)\longrightarrow i_s^*(\mathcal O_X)=\mathcal O_{X_s}.\]
\begin{SGAresult}{命题}{2.2}
\footnote{由 A. Douady 告知编者。}
设 $X,Y$ 是 $S$-\Term{mixed-manifold}{混合流形}，$s\in S$。假设\Term{fiber}{纤维} $X_s$ \Term{paracompact}{仿紧}，且 $s$ 有一个由\Term{paracompact}{仿紧}\Term{neighborhood}{邻域}组成的\Term{neighbourhood-basis}{基本邻域系}。设 $A$ 是 $X_s$ 的\Term{closed-subset}{闭集}，$f\in H^0(A,\mathcal Hom(X_s,Y_s))$。令 $F$ 为 $i_s^{-1}\mathcal Hom(X,Y)$ 的\Term{subsheaf}{子层}，其芽由那些“延拓 $f$ 的芽”组成（即其经 2.1.1 的像在 $A$ 上与 $f$ 一致）。则 $F$ 是\Term{soft-sheaf}{软层}。
\end{SGAresult}
\begin{proof}
设 $X'$ 是 $x$ 在 $X$ 中的\Term{open-neighborhood}{开邻域}，作为\Term{mixed-manifold}{混合流形}\Term{isomorphic}{同构}于 $S'\times U$，其中 $S'$ 是 $s$ 的\Term{paracompact}{仿紧}\Term{neighborhood}{邻域}，$U$ 是 $\mathbb R^m$ 中\Term{relatively-compact}{相对紧}的\Term{open-subset}{开集}，并且 $X'_s\cap A$ 经 $f$ 的像包含在 $Y$ 的一个\Term{open-subset}{开集}内，该\Term{open-subset}{开集}\Term{isomorphic}{同构}于 $S'\times V$，其中 $V$ 是 $\mathbb R^n$ 中包含原点的凸\Term{open-subset}{开集}。由 T.F. II 3.4.1，只需证明 $F$ 在 $X'_s$ 中的\Term{closed-subset}{闭集} $A_0\subset X_s$ 上的每个\Term{section}{截面}都能延拓到 $X'_s$。而 $X'_s$（$\simeq\{s\}\times U$）是\Term{relatively-compact}{相对紧}的，所以 $X_s$ 中包含在 $X'_s$ 内的\Term{closed-subset}{闭集}是紧的。将 $X'$（相应地 $Y'$）等同于 $S'\times U$（相应地 $S'\times V$），令 $A'=A\cap X'$，并令 $A_0$ 为 $X'_s=\{s\}\times U$ 中的紧集。由于 $A'$ 是\Term{paracompact}{仿紧}\Term{space}{空间} $\{s\}\times U$ 中的\Term{closed-subset}{闭集}，由 T.F. II 3.3.1\footnote{\Term{paracompact}{仿紧}\Term{space}{空间}的\Term{closed-subset}{闭集}有一个由\Term{paracompact}{仿紧}\Term{neighborhood}{邻域}组成的\Term{neighbourhood-basis}{基本邻域系}。}，有
\[H^0(A',\mathcal Hom(X_s,Y_s))=\varinjlim_{U'\supset A'}\operatorname{Hom}(U',Y'_s),\]
\Term{limit}{极限}沿 $X'_s$ 中包含 $A'$ 的\Term{open-subset}{开集} $U'$ 取。类似地，$A_0$ 是\Term{paracompact}{仿紧}\Term{space}{空间} $S'\times\overline U$ 中的\Term{closed-subset}{闭集}（\Term{paracompact}{仿紧}\Term{space}{空间}与\Term{compact-space}{紧空间}的积仍\Term{paracompact}{仿紧}），所以
\[H^0(A_0,i_s^{-1}\mathcal Hom(X,Y))=\varinjlim_{U_0\supset A_0}\operatorname{Hom}(U_0,Y'),\]
\Term{limit}{极限}沿 $X'$ 中包含 $A_0$ 的\Term{open-subset}{开集} $U_0$ 取。因此命题将由下列\Term{lemma}{引理}\Term{pushout}{推出}。
\end{proof}
\SourcePage{216}
\begin{SGAresult}{\Term{lemma}{引理}}{2.2.1}
设 $U$ 是 $\mathbb R^m$ 的\Term{open-subset}{开集}，$V$ 是 $\mathbb R^n$ 中包含 $0$ 的凸\Term{open-subset}{开集}。设 $A_0,A_1$ 是 $\{s\}\times U$ 的\Term{closed-subset}{闭集}，$U_0$（相应地 $U_1$）是 $S\times U$（相应地 $\{s\}\times U$）中包含 $A_0$（相应地 $A_1$）的\Term{open-subset}{开集}。设 $f_1:U_1\to V$ 是一个\Term{morphism}{态射}，$f_0:U_0\to S\times V$ 是一个 $S$-\Term{morphism}{态射}，它在 $U_0\cap(\{s\}\times U)$ 上诱导的\Term{morphism}{态射}在 $A_1$ 附近定义与 $f_1$ 相同的芽。那么存在 $S\times U$ 中包含 $s\times U$ 的\Term{open-subset}{开集} $U'$，以及 $S$-\Term{morphism}{态射} $\widetilde f:U'\to S\times V$，使得：
\begin{enumerate}
\item[(i)] $\widetilde f$ 沿 $\{s\}\times U$ 的芽延拓 $f_0$ 定义的芽；
\item[(ii)] $\widetilde f$ 在 $\{s\}\times U$ 上诱导的\Term{morphism}{态射}在 $A_1$ 附近定义与 $f_1$ 相同的芽。
\end{enumerate}
\end{SGAresult}
\begin{proof}
以 $af_1$ 代替 $f_1$，其中 $a$ 是 $U$ 上的 $C^r$ 函数，在 $A_1$ 附近等于 $1$ 且\Term{support}{支集}包含在 $U_1$ 中，即可假定 $U_1=U$。设 $b$ 是 $U$ 上的 $C^r$ 函数，在 $A_0$ 上等于 $1$ 且\Term{support}{支集}包含在 $U_0\cap(\{s\}\times U)$ 中。那么
\[U'=U_0\cup\bigl(S\times(U-\operatorname{Supp}b)\bigr)\]
是 $S\times U$ 中包含 $\{s\}\times U$ 的\Term{open-subset}{开集}。由
\[\widetilde f(t,x)=\bigl(t,b(x)\operatorname{pr}_2f_0(t,x)+(1-b(x))f_1(x)\bigr)\]
定义的\Term{morphism}{态射} $\widetilde f:U'\to S\times V$ 显然满足 (i)、(ii)。
\end{proof}
\begin{SGAresult}{推论}{2.3}
在 2.2 的假设下，\Term{sheaf}{层} $i_s^{-1}\mathcal Hom(X,Y)$、$\mathcal Hom(X_s,Y_s)$ 是\Term{soft-sheaf}{软层}。特别地，\Term{sheaf}{层} $i_s^{-1}(\mathcal O_X)$、$\mathcal O_{X_s}$ 是\Term{soft-sheaf}{软层}\footnote{对于 $\mathcal O_{X_s}$，这“众所周知”。编者不知道：若对每个 $s\in S$ 均满足 2.2 的假设且 $X$ \Term{paracompact}{仿紧}，$\mathcal O_X$ 是否为\Term{soft-sheaf}{软层}。借助上述假设，这将显然\Term{pushout}{推出} $i_s^{-1}(\mathcal O_X)$ 是\Term{soft-sheaf}{软层}以及下面的 2.5。}。对 $X_s$ 的每个\Term{closed-subset}{闭集} $A$，由 2.1.1 得到的映射
\[H^0(X_s,i_s^{-1}\mathcal Hom(X,Y))\longrightarrow H^0(A,\mathcal Hom(X_s,Y_s))\]
是\Term{surjective}{满射}。
\end{SGAresult}
\SourcePage{217}
\begin{proof}
对 $A=\varnothing$ 应用 2.2 即知 $i_s^{-1}\mathcal Hom(X,Y)$ 是\Term{soft-sheaf}{软层}；取 $S=s$，可知 $\mathcal Hom(X_s,Y_s)$ 是\Term{soft-sheaf}{软层}。最后一个断言由 2.2 得出。
\end{proof}
\begin{SGAresult}{推论}{2.4}
在 2.2 的假设下，设 $E$ 是一个 $\mathcal O_X$-\Term{module}{模}。对 $X_s$ 的每个\Term{closed-subset}{闭集} $A$，由\Term{epimorphism}{满射}
\[i_s^{-1}(E)\longrightarrow i_s^*(E)=i_s^{-1}(E)\otimes_{i_s^{-1}(\mathcal O_X)}\mathcal O_{X_s}\]
得到的\Term{canonical}{典范}\Term{homomorphism}{同态}
\[H^0(X_s,i_s^{-1}(E))\longrightarrow H^0(A,i_s^*(E))\]
是\Term{epimorphism}{满射}。
\end{SGAresult}
\begin{proof}
由于 $i_s^{-1}(\mathcal O_X)$ 是\Term{soft-sheaf}{软层}（2.3），\Term{functor}{函子} $H^0(A,\ )$ 在 $i_s^{-1}(\mathcal O_X)$-\Term{module}{模}的\Term{category}{范畴}上是\Term{exact-functor}{正合}的（T.F. II 3.7.1、3.4.2、4.4.3），所以\Term{homomorphism}{同态}
\[H^0(A,i_s^{-1}(E))\longrightarrow H^0(A,i_s^*(E))\]
是\Term{epimorphism}{满射}。由于 $i_s^{-1}(E)$ 是\Term{soft-sheaf}{软层}，\Term{homomorphism}{同态}
\[H^0(X_s,i_s^{-1}(E))\longrightarrow H^0(A,i_s^{-1}(E))\]
也是\Term{epimorphism}{满射}，于是得证。
\end{proof}
\begin{SGAresult}{推论}{2.5}
设 $p:X\to S$ 是一个\Term{mixed-manifold}{混合流形}，$s\in S$。假设 $s$ 有一个由\Term{paracompact}{仿紧}\Term{neighborhood}{邻域} $S'$ 组成的\Term{neighbourhood-basis}{基本邻域系}，使各 $p^{-1}(S')$ 构成 $X_s$ 的\Term{paracompact}{仿紧}\Term{neighbourhood-basis}{基本邻域系}（例如当 $S$ \Term{locally-compact}{局部紧}而 $p$ 是\Term{topological-space}{拓扑空间}之间的\Term{proper-map}{固有映射}时，这个条件成立）。那么对每个 $\mathcal O_X$-\Term{module}{模} $E$，有
\[\mathrm R^ip_*(E)_s=0\qquad(i>0).\]
\end{SGAresult}
\SourcePage{218}
\begin{proof}
由 T.F. II 4.17.1，有
\[\mathrm R^ip_*(E)_s=H^i(X_s,i_s^{-1}(E)).\]
但由于 $i_s^{-1}(\mathcal O_X)$ 是\Term{soft-sheaf}{软层}（2.3），有
\[H^i(X_s,i_s^{-1}(E))=0\qquad(i>0),\]
断言随之得出。
\end{proof}
\begin{SGAresult}{推论}{2.6}
设 $X,Y$ 是 $S$-\Term{mixed-manifold}{混合流形}，$s\in S$。假设 $X$ 满足 2.5 的假设。对每个\Term{morphism}{态射} $f:X_s\to Y_s$，存在 $s$ 的\Term{open-neighborhood}{开邻域} $S'$ 及\Term{morphism}{态射} $\overline f:X|S'\to Y|S'$，使 $\overline f_s=f$。
\end{SGAresult}
\begin{proof}
这直接由 2.3（取 $A=X_s$）以及在 $X$ 的假设下成立的等式
\[H^0(X_s,i_s^{-1}\mathcal Hom(X,Y))=\varinjlim_{S'\ni s}\operatorname{Hom}(X|S',Y)\qquad(S'\text{ 开})\]
\Term{pushout}{推出}（T.F. II 3.3.1）。
\end{proof}
\begin{SGAresult}{命题}{2.7}
设 $X,Y$ 是 $S$-\Term{mixed-manifold}{混合流形}，$f:X\to Y$ 是\Term{morphism}{态射}，$s\in S$，$x\in X_s$。若 $f_s:X_s\to Y_s$ 在 $x$ 处的切\Term{linear}{线性}映射\Term{invertible}{可逆}，则 $f$ 在 $x$ 处是\Term{local-isomorphism}{局部同构}，即存在 $x$ 在 $X$ 中的\Term{open-neighborhood}{开邻域} $U$ 以及 $Y$ 中的\Term{open-subset}{开集} $V$，使 $f$ 诱导 $U$ 到 $V$ 的\Term{isomorphism}{同构}。
\end{SGAresult}
\begin{proof}
这是\Term{implicit-function-theorem}{隐函数定理}的简单推论。细节留给读者。
\end{proof}
\begin{SGAresult}{推论}{2.8}
设 $X,Y$ 是 $S$-\Term{mixed-manifold}{混合流形}，$s\in S$。假设 $s$ 有一个开\Term{neighbourhood-basis}{基本邻域系} $S'$，使各 $X|S'$（相应地 $Y|S'$）构成 $X_s$（相应地 $Y_s$）的\Term{neighbourhood-basis}{基本邻域系}（例如 $X$（相应地 $Y$）在 $S$ 上\Term{proper}{固有}时条件成立）。设 $f:X\to Y$ 是\Term{morphism}{态射}，且 $f_s:X_s\to Y_s$ 是\Term{isomorphism}{同构}。那么存在 $s$ 的\Term{neighborhood}{邻域} $S'$，
\SourcePage{219}
使 $f|X|S':X|S'\to Y|S'$ 是\Term{isomorphism}{同构}。
\end{SGAresult}
\begin{proof}
由 2.7，存在 $X_s$ 在 $X$ 中的\Term{open-neighborhood}{开邻域} $U$，使 $f|U$ \Term{etale}{平展}（即为\Term{local-isomorphism}{局部同构}）。可以假定 $U=X|S'$，其中 $S'$ 是 $s$ 在 $S$ 中的\Term{neighborhood}{邻域}；换言之，缩小 $S$ 后可假定 $f$ \Term{etale}{平展}。于是容易验证，\Term{fiber-product}{纤维积} $Z=X\times_YX$ 自然具有 $S$-\Term{mixed-manifold}{混合流形}结构（各\Term{projection}{投影}均\Term{etale}{平展}于 $X$），而\Term{diagonal-morphism}{对角态射} $\Delta:X\to Z$ 是\Term{open-immersion}{开浸入}\footnote{注意与\Term{algebraic-geometry}{代数几何}的类比（SGA 1 I 5.1）。}。另一方面立即可见，$Z$ 与 $X,Y$ 满足相同假设，即 $S'$ 遍历 $s$ 在 $S$ 中的\Term{open-neighborhood}{开邻域}时，$Z|S'$ 型的\Term{open-subset}{开集}构成 $Z_s$ 的\Term{neighbourhood-basis}{基本邻域系}。由于 $\Delta_s:X_s\to Z_s$ 是\Term{isomorphism}{同构}，进一步缩小 $S$ 后，$\Delta$ 就是\Term{isomorphism}{同构}。因此 $f$ 是\Term{open-immersion}{开浸入}；再缩小一次 $S$，它就是\Term{isomorphism}{同构}，于是得证。
\end{proof}
\begin{SGAresult}{推论}{2.9}
设 $X$ 是一个 $S$-\Term{mixed-manifold}{混合流形}，$s\in S$。假设 2.5 的条件成立。那么存在 $s$ 在 $S$ 中的\Term{neighborhood}{邻域} $S'$，使 $X|S'$ 与 $S'\times X_s$ 在 $S'$ 上\Term{isomorphic}{同构}。
\end{SGAresult}
\begin{proof}应用 2.6 及 2.8。\end{proof}
\begin{SGAresult}{推论}{2.10}
设 $s\in S$，$V$ 是\Term{paracompact}{仿紧}\Term{manifold}{流形}。置 $X=S\times V$，以 $i_s$ 记由 $s$ 定义的 $V$ 到 $X$ 的\Term{embedding}{嵌入}。假设 2.5 的条件成立。设 $E$ 是\Term{locally-free-finite-type}{有限型局部自由} $\mathcal O_X$-\Term{module}{模}（相应地 $(\mathcal O_X\otimes\mathbb C)$-\Term{module}{模}）。那么存在 $s$ 在 $S$ 中的\Term{neighborhood}{邻域} $S'$，使 $E|X|S'$ \Term{isomorphic}{同构}于 $\operatorname{pr}_2^*i_s^*(E)$（$\operatorname{pr}_2:S'\times V\to V$ 是第二\Term{projection}{投影}）。
\end{SGAresult}
\begin{proof}
置 $F=\operatorname{pr}_2^*i_s^*(E)$。对\Term{module}{模} $\mathcal Hom(E,F)$ 应用 2.4，可知诱导\Term{homomorphism}{同态}
\[\operatorname{Hom}(i_s^{-1}(E),i_s^{-1}(F))\longrightarrow\operatorname{Hom}(i_s^*(E),i_s^*(F))\]
\SourcePage{220}
是\Term{epimorphism}{满射}。另一方面，由 T.F. II 3.3.1，有
\begin{align*}
\operatorname{Hom}(i_s^{-1}(E),i_s^{-1}(F))
&\simeq H^0(X_s,i_s^{-1}\mathcal Hom(E,F))\\
&=\varinjlim_{S'\ni s}H^0(X|S',\mathcal Hom(E,F))\quad(S'\text{ 开})\\
&\simeq\varinjlim_{S'\ni s}\operatorname{Hom}(E|X|S',F|X|S').
\end{align*}
因此\Term{canonical}{典范}\Term{isomorphism}{同构} $i_s^*(E)\simeq i_s^*(F)$ 由一个\Term{homomorphism}{同态} $f:E\to F$ 诱导（必要时缩小 $S$）。由于 $f_s$ 是\Term{isomorphism}{同构}，由 \Term{nakayama-lemma}{Nakayama 引理}，$f$ 在 $X_s$ 附近是\Term{isomorphism}{同构}，从而在 $X|S'$ 上是\Term{isomorphism}{同构}，其中 $S'$ 为 $s$ 的适当\Term{neighborhood}{邻域}，证毕。
\end{proof}

\SGAsection{3}{有限性定理}
\textbf{3.0. 约定及记号。}
\textbf{3.0.1.} 设 $X,Y$ 是\Term{topological-space}{拓扑空间}。以 $C(X,Y)$ 记从 $X$ 到 $Y$ 的\Term{continuous-map}{连续映射}\Term{set}{集合}，以 $\mathcal C(X,Y)$ 记 $X$ 上由 $U\mapsto C(U,Y)$ 定义的\Term{sheaf}{层}。

\textbf{3.0.2.} 设 $E,F$ 是实或复 \Term{banach-space}{Banach 空间}。以 $L(E,F)$ 记 $E$ 到 $F$ 的\Term{continuous}{连续}\Term{linear}{线性}映射所成的 \Term{banach-space}{Banach 空间}。

\textbf{3.0.3.} 本节讨论的\Term{mixed-manifold}{混合流形}均假定为 $C^\infty$ 类。设 $X$ 是一个 $S$-\Term{mixed-manifold}{混合流形}，我们经常把\Term{locally-free-finite-type}{有限型局部自由} $\mathcal O_X$-\Term{module}{模}（相应地 $(\mathcal O_X\otimes\mathbb C)$-\Term{module}{模}）与相应的实（相应地复）\Term{vector-bundle}{向量丛}等同。

\textbf{3.0.4.} 设 $X$ 是一个通常意义下的紧\Term{manifold}{流形}，$E$ 是 $X$ 上\Term{finite-rank}{有限秩}实或复\Term{vector-bundle}{向量丛}。以 $C^\infty(X,E)$ 记 $E$ 在 $X$ 上的 $C^\infty$ \Term{section}{截面}的 Fréchet \Term{space}{空间}。对 $r\in\mathbb Z$，以 $C_2^r(X,E)$ 记 $E$ 的 $L^2$ 意义下 $r$ 次可微\Term{section}{截面}的 Hilbert \Term{space}{空间}（[16] IX）。对 $s\geq r$，$C_2^s(X,E)$ 通过一个\Term{continuous}{连续}的\Term{injective-map}{单射}映到 $C_2^r(X,E)$，其像\Term{dense}{稠密}，而且
\SourcePage{221}
有（同上）
\[C^\infty(X,E)=\varprojlim_r C_2^r(X,E).\]
\textbf{3.0.5.} 设 $S$ 是\Term{topological-space}{拓扑空间}。我们需要 $S$ 上的“\Term{hilbert-bundle}{Hilbert 丛}”概念，读者可参阅 [12] 或 [14]。有一个从 $S$ 上实（相应地复）\Term{hilbert-bundle}{Hilbert 丛}的\Term{category}{范畴}到 $\mathcal O_S$-\Term{module}{模}（相应地 $(\mathcal O_S\otimes\mathbb C)$-\Term{module}{模}）\Term{category}{范畴}的\Term{additive-functor}{加性函子}，它把每个\Term{hilbert-bundle}{Hilbert 丛} $E$ 送到其\Term{continuous}{连续}\Term{section}{截面}\Term{sheaf}{层} $\mathcal E$。

\textbf{3.1. 辅助技术结果。}
\begin{SGAresult}{\Term{lemma}{引理}}{3.1.1}
设 $X,Y$ 是\Term{topological-space}{拓扑空间}，$X$ \Term{locally-compact}{局部紧}。以 $\mathcal C'(X,Y)$ 记在 $X$ 的\Term{relatively-compact}{相对紧}\Term{open-basis}{开集基}上由 $\mathcal C'(X,Y)(U)=C(\overline U,Y)$ 定义的\Term{presheaf}{预层}（3.0.1）。\Term{canonical}{典范}\Term{restriction}{限制}\Term{morphism}{态射}
\[\mathcal C'(X,Y)\longrightarrow\mathcal C(X,Y)\]
在对应的\Term{associated-sheaf}{伴随层}上诱导\Term{isomorphism}{同构}。
\end{SGAresult}
\begin{proof}这是直接的。\end{proof}
\begin{SGAresult}{\Term{lemma}{引理}}{3.1.2}
设 $E,F$ 是\Term{locally-convex}{局部凸}\Term{topology}{拓扑}\Term{vector-space}{向量空间}（为明确起见，设它们在 $\mathbb C$ 上），$u:E\to F$ 是\Term{continuous}{连续}\Term{linear}{线性}映射（相应地\Term{continuous}{连续}且像\Term{dense}{稠密}的\Term{linear}{线性}映射）。设 $X$ 是\Term{compact-space}{紧空间}。在 $C(X,E),C(X,F)$ 上赋予一致收敛\Term{topology}{拓扑}，则由 $u$ 诱导的\Term{linear}{线性}映射
\[u_X:C(X,E)\longrightarrow C(X,F)\]
\Term{continuous}{连续}（相应地\Term{continuous}{连续}且像\Term{dense}{稠密}）。
\end{SGAresult}
\SourcePage{222}
\begin{proof}
显然 $u_X$ \Term{continuous}{连续}。另一方面（Bourbaki，《积分》，第 III 章第 2 节，\Term{lemma}{引理} 1、2），$C(X,\mathbb C)\otimes_{\mathbb C}E$ 在 $C(X,E)$ 中\Term{dense}{稠密}（它由那些取值于一个\Term{finite-dimensional}{有限维}\Term{vector-subspace}{向量子空间}的\Term{continuous-map}{连续映射} $X\to E$ 组成）；类似地，$C(X,\mathbb C)\otimes_{\mathbb C}F$ 在 $C(X,F)$ 中\Term{dense}{稠密}。假设 $u$ \Term{continuous}{连续}且像\Term{dense}{稠密}。那么
\[u_X\bigl(C(X,\mathbb C)\otimes_{\mathbb C}E\bigr)=C(X,\mathbb C)\otimes_{\mathbb C}u(E)\]
在 $C(X,\mathbb C)\otimes_{\mathbb C}F$ 中\Term{dense}{稠密}。取\Term{closure}{闭包}即知 $u_X(C(X,E))$ 在 $C(X,F)$ 中\Term{dense}{稠密}，证毕。
\end{proof}
\textbf{3.1.3.} 设 $V$ 是紧\Term{manifold}{流形}，$E,F$ 是 $V$ 上的\Term{finite-rank}{有限秩}复\Term{vector-bundle}{向量丛}\footnote{为明确起见如此设定；实数情形有类似结果。}，$n\in\mathbb N$。由 1.2.3.1，\Term{space}{空间} $\operatorname{Diff}_{\mathbb C}^n(E,F)$ 等同于
\[\operatorname{Hom}(\mathcal P^n(E),F)\simeq C^\infty(V,\mathcal Hom(\mathcal P^n(E),F)),\]
因而自然具有 Fréchet \Term{space}{空间}结构。设 $u\in\operatorname{Diff}_{\mathbb C}^n(E,F)$；容易验证，由 $u$ 定义的从 $C^\infty(V,E)$ 到 $C^\infty(V,F)$ 的\Term{linear}{线性}映射\Term{continuous}{连续}。若 $L(C^\infty(V,E),C^\infty(V,F))$ 表示从 $C^\infty(V,E)$ 到 $C^\infty(V,F)$ 的\Term{continuous}{连续}\Term{linear}{线性}映射\Term{space}{空间}，则有\Term{linear}{线性}映射
\[\phi:\operatorname{Diff}_{\mathbb C}^n(E,F)\longrightarrow L(C^\infty(V,E),C^\infty(V,F)).\]
另一方面，对 $r\in\mathbb Z$，$u$ 通过\Term{scalar-extension}{标量扩张}（[16] IX \Term{theorem}{定理} 12）\Term{action}{作用}于 $E$ 的 $C_2^r$ \Term{section}{截面}，并定义\Term{linear}{线性}映射
\[\phi_r:\operatorname{Diff}_{\mathbb C}^n(E,F)\longrightarrow L(C_2^r(V,E),C_2^{r-n}(V,F)).\]
容易验证：
\begin{enumerate}
\item[(i)] 各 $\phi_r$ \Term{continuous}{连续}（目标\Term{space}{空间}是 \Term{banach-space}{Banach 空间}，3.0.2）；
\SourcePage{223}
\item[(ii)] $\phi=\varprojlim_r\phi_r$（即对 $u\in\operatorname{Diff}_{\mathbb C}^n(E,F)$，$\phi(u)$ 是由各 $\phi_r(u)$ 组成的\Term{projective-system}{投射系}的\Term{limit}{极限}）。
\end{enumerate}
上述断言由所引文献的结果很容易得到：立即归结到 $n=0$；此时 (ii) 直接由定义得出，而 (i) 则由乘积的高阶导数公式得到（先借助\Term{partition-of-unity}{单位分解}将问题\Term{localization}{局部化}）。

\begin{SGAresult}{\Term{lemma}{引理}}{3.1.4}
设 $S$ 是\Term{topological-space}{拓扑空间}，$V$ 是紧\Term{manifold}{流形}，$E,F$ 是 $V$ 上的\Term{finite-rank}{有限秩}复\Term{vector-bundle}{向量丛}，$n\in\mathbb N$。以 $X$ 记 $S$-\Term{mixed-manifold}{混合流形} $S\times V$。对 $u\in\operatorname{Diff}_{\mathbb C}^n(\operatorname{pr}_2^*(E),\operatorname{pr}_2^*(F))$，从 $S$ 到 $\operatorname{Diff}_{\mathbb C}^n(E,F)$ 的映射 $s\mapsto i_s^*(u)$ \Term{continuous}{连续}。由此定义的\Term{linear}{线性}映射
\[\operatorname{Diff}_{\mathbb C}^n(\operatorname{pr}_2^*(E),\operatorname{pr}_2^*(F))\longrightarrow C(S,\operatorname{Diff}_{\mathbb C}^n(E,F))\]
是\Term{isomorphism}{同构}。
\end{SGAresult}
\begin{proof}
利用 1.2.3.1 的等同，只需考虑 $n=0$。令 $\mathcal Hom(E,F)=G$，问题就归结为证明：对 $u\in\Gamma(X,\operatorname{pr}_2^*(G))$，从 $S$ 到 $C^\infty(V,G)$ 的映射 $s\mapsto u_s=i_s^*(u)$ \Term{continuous}{连续}，且由此得到的\Term{linear}{线性}映射
\[\Gamma(X,\operatorname{pr}_2^*(G))\longrightarrow C(S,C^\infty(V,G))\]
是\Term{isomorphism}{同构}。但这直接由\Term{mixed-manifold}{混合流形}的定义（1.1）得到。
\end{proof}
\textbf{3.1.5.} 在 3.1.4 的假设下，设
\[u\in\operatorname{Diff}^n(\operatorname{pr}_2^*(E),\operatorname{pr}_2^*(F)).\footnote{为简便起见，省去下标 $\mathbb C$。}\]
\SourcePage{224}
由 3.1.4，将 $u$ 等同于从 $S$ 到 $\operatorname{Diff}^n(E,F)$ 的\Term{continuous-map}{连续映射}。可见，对 $S$ 中的\Term{open-subset}{开集} $U$，\Term{morphism}{态射}
\[\operatorname{pr}_{1*}(u)(U):\Gamma(X|U,E)\longrightarrow\Gamma(X|U,F)\]
\Term{canonical}{典范}地等同于由
\[\phi u:U\longrightarrow L(C^\infty(V,E),C^\infty(V,F))\qquad(3.1.3)\]
得到的\Term{linear}{线性}映射
\[C(U,C^\infty(V,E))\longrightarrow C(U,C^\infty(V,F)).\]
假设 $S$ \Term{locally-compact}{局部紧}。由 3.1.1，\Term{morphism}{态射}
\[\operatorname{pr}_{1*}(u):\operatorname{pr}_{1*}\operatorname{pr}_2^*(E)\longrightarrow\operatorname{pr}_{1*}\operatorname{pr}_2^*(F)\]
就是由 $\phi u$ 定义的\Term{presheaf}{预层}\Term{morphism}{态射}
\[u_\infty:\mathcal C'(S,C^\infty(V,E))\longrightarrow\mathcal C'(S,C^\infty(V,F))\]
所对应的\Term{sheaf}{层}\Term{morphism}{态射}。另一方面，对 $r\in\mathbb Z$，
\[\phi_ru:S\longrightarrow L(C_2^r(V,E),C_2^{r-n}(V,F))\]
定义\Term{presheaf}{预层}\Term{morphism}{态射}
\[u_r:\mathcal C'(S,C_2^r(V,E))\longrightarrow\mathcal C'(S,C_2^{r-n}(V,F)).\]
由 3.1.3(ii)，各 $u_r$ 构成\Term{projective-system}{投射系}，\Term{limit}{极限}为 $u_\infty$。将 $\phi_ru$ 解释为\Term{trivial-hilbert-bundle}{平凡 Hilbert 丛}之间的\Term{morphism}{态射}（3.0.5）对我们很方便：
\[S\times C_2^r(V,E)\longrightarrow S\times C_2^{r-n}(V,F),\]
\SourcePage{225}
它在\Term{section}{截面}\Term{sheaf}{层}上诱导与 $u_r$ 对应的\Term{sheaf}{层}\Term{morphism}{态射}。

\textbf{3.2. 指标的定义。}
\begin{SGAresult}{\Term{theorem}{定理}}{3.2.1}
设 $S$ 是\Term{locally-compact-space}{局部紧空间}，$p:X\to S$ 是\Term{proper}{固有}\Term{mixed-manifold}{混合流形}，$E^\bullet$ 是 $X$ 上复\Term{vector-bundle}{向量丛}组成的阶数为 $\underline n$ 的\Term{elliptic-complex}{椭圆复形}（1.2.5）。则\Term{complex}{复形} $\mathrm Rp_*(E^\bullet)$ 是\Term{perfect-complex}{完美复形}（由 2.5，它在 $(\mathcal O_S\otimes\mathbb C)$-\Term{module}{模}\Term{category}{范畴}的\Term{derived-category}{导出范畴}中\Term{isomorphic}{同构}于 $p_*(E^\bullet)$）。
\end{SGAresult}
\begin{proof}
问题在 $S$ 上是\Term{local-scheme}{局部}的，所以由 2.9 可假定 $X=S\times V$，其中 $V$ 是紧\Term{manifold}{流形}，且 $p=\operatorname{pr}_1$。此外，由 2.10 可假定对每个 $i$ 有 $E^i=\operatorname{pr}_2^*(F^i)$，其中 $F^i$ 是 $V$ 上\Term{finite-rank}{有限秩}复\Term{vector-bundle}{向量丛}。由 3.1.5，\Term{complex}{复形} $p_*(E)$ 对应于\Term{presheaf}{预层}\Term{complex}{复形} $(\mathcal C'(S,C^\infty(V,F^i)),d_\infty^i)$（其中 $d^i:E^i\to E^{i+1}$ 是给定的\Term{differential-operator}{微分算子}）；简记这个\Term{complex}{复形}为 $\mathcal C'(S,C^\infty(V,F^\bullet))$。对每个 $i\in\mathbb Z$，置
\[m_i=\sum_{j<i}n_j\qquad(\text{约定 }E^i=0\text{ 时 }n_i=0).\]
对每个 $r\in\mathbb Z$，由 3.1.5 得到\Term{presheaf}{预层}\Term{complex}{复形} $(\mathcal C'(S,C_2^{r-m_i}(V,F^i)),d^i_{r-m_i})$，记为 $\mathcal C'(S,C_2^r(V,F^\bullet))$。各 $\mathcal C'(S,C_2^r(V,F^\bullet))$ 构成\Term{projective-system}{投射系}，且由 3.1.5 有
\[\mathcal C'(S,C^\infty(V,F^\bullet))=\varprojlim_r\mathcal C'(S,C_2^r(V,F^\bullet)).\tag{3.2.1.1}\]
最后，以 $F_r^\bullet$ 记 $S$ 上由各 $d^i_{r-m_i}$ 定义的\Term{trivial-hilbert-bundle}{平凡 Hilbert 丛}\Term{complex}{复形}，以 $\mathcal F_r^\bullet$ 记相应\Term{section}{截面}\Term{sheaf}{层}的\Term{complex}{复形}；由 3.1.5，它正是与 $\mathcal C'(S,C_2^r(V,F^\bullet))$ 对应的\Term{sheaf}{层}\Term{complex}{复形}。至此，为证明\Term{theorem}{定理}，只需证明下列断言：

(a) 对 $r\in\mathbb Z$，\Term{canonical}{典范}\Term{morphism}{态射}
\SourcePage{226}
\[\mathcal C'(S,C^\infty(V,F^\bullet))\longrightarrow\mathcal C'(S,C_2^r(V,F^\bullet))\]
是\Term{presheaf}{预层}\Term{complex}{复形}的\Term{quasi-isomorphism}{拟同构}，即在\Term{cohomology-presheaf}{上同调预层}上诱导\Term{isomorphism}{同构}。

(b) 对 $r\in\mathbb Z$，\Term{complex}{复形} $\mathcal F_r^\bullet$ 是\Term{perfect-complex}{完美复形}。

\textbf{(b) 的证明。} $E^\bullet$ 的椭圆性蕴含（[16] XI \Term{theorem}{定理} 6\footnote{它容易推广到\Term{elliptic-complex}{椭圆复形}。}）：对每个 $s\in S$，$(F_r^\bullet)_s$ 的\Term{differential}{微分}具有\Term{scheme-theoretic-image}{闭像}，且各 $H^i((F_r^\bullet)_s)$ 是\Term{finite-dimensional-vector-space}{有限维向量空间}。换言之，采用 [11] 的术语，$F_r^\bullet$ 是\Term{hilbert-bundle}{Hilbert 丛}的一个\Term{direct-quasi-acyclic-complex}{直接拟无环复形}（complexe quasi-acyclique direct）。因此由 [12] 2.3，$F_r^\bullet$ \Term{local-scheme}{局部}\Term{homotopy-equivalent}{同伦等价}于一个\Term{finite-rank}{有限秩}\Term{vector-bundle}{向量丛}的\Term{bounded}{有界}\Term{complex}{复形}。由于“\Term{section}{截面}\Term{sheaf}{层}”\Term{functor}{函子}是\Term{additive-functor}{加性}的（3.0.5），$\mathcal F_r^\bullet$ \Term{local-scheme}{局部}\Term{homotopy-equivalent}{同伦等价}于一个\Term{locally-free-finite-type}{有限型局部自由}\Term{sheaf}{层}的\Term{bounded}{有界}\Term{complex}{复形}，所以它是\Term{perfect-complex}{完美复形}。

为证明 (a)，还需下列中间断言：

(a') 对 $p\geq q$，\Term{canonical}{典范}\Term{morphism}{态射}
\[\mathcal F_p^\bullet\longrightarrow\mathcal F_q^\bullet\]
是\Term{quasi-isomorphism}{拟同构}。

\textbf{(a') 的证明。} 以 $M_{pq}^\bullet$ 记\Term{morphism}{态射} $F_p^\bullet\to F_q^\bullet$ 的\Term{mapping-cone}{锥}这一\Term{hilbert-bundle}{Hilbert 丛}\Term{complex}{复形}。$E$ 的椭圆性蕴含：对每个 $s\in S$，\Term{morphism}{态射} $(F_p^\bullet)_s\to(F_q^\bullet)_s$ 是\Term{quasi-isomorphism}{拟同构}（[16] XI \Term{theorem}{定理} 6 的推论 2）。换言之，$M_{pq}^\bullet$ 是\Term{direct-acyclic-complex}{直接无环复形}。由 [11] 2.3 可知 $M_{pq}^\bullet$ \Term{local-scheme}{局部}\Term{homotopically-trivial}{同伦平凡}，所以 $\mathcal M_{pq}^\bullet$ 是\Term{acyclic}{无上同调}的；而它正是 $\mathcal F_p^\bullet\to\mathcal F_q^\bullet$ 的\Term{mapping-cone}{锥}，证毕。

\textbf{(a) 的证明。} 需证明，对 $S$ 的每个\Term{relatively-compact}{相对紧}\Term{open-subset}{开集} $U$，\Term{morphism}{态射}
\SourcePage{227}
\[C(\overline U,C^\infty(V,F^\bullet))\longrightarrow C(\overline U,C_2^r(V,F^\bullet))\]
是\Term{quasi-isomorphism}{拟同构}。将 $S$ 换成 $\overline U$（并将 $E^\bullet$ 换成在 $\overline U$ 上诱导的\Term{complex}{复形}），问题就归结为：在假定 $S$ 紧时证明
\[C(S,C^\infty(V,F^\bullet))\longrightarrow C(S,C_2^r(V,F^\bullet))\]
是\Term{quasi-isomorphism}{拟同构}。由 3.2.1.1 有
\[C(S,C^\infty(V,F^\bullet))=\varprojlim_r C(S,C_2^r(V,F^\bullet)).\]
对每个 $i$，$C(S,C_2^r(V,F^i))$ 是 \Term{banach-space}{Banach 空间}（赋予一致收敛\Term{topology}{拓扑}）；由 3.0.4 及 3.1.2，过渡\Term{morphism}{态射}
\[C(S,C_2^p(V,F^i))\longrightarrow C(S,C_2^q(V,F^i))\]
是像\Term{dense}{稠密}的\Term{continuous}{连续}\Term{linear}{线性}映射。另一方面，对 (a') 的\Term{quasi-isomorphism}{拟同构} $\mathcal F_p^\bullet\to\mathcal F_q^\bullet$ 应用\Term{functor}{函子} $\Gamma_S$（因 $\mathcal O_S$ 是\Term{soft-sheaf}{软层}，它是\Term{exact-functor}{正合}的），可知过渡\Term{morphism}{态射}
\[C(S,C_2^p(V,F^\bullet))\longrightarrow C(S,C_2^q(V,F^\bullet))\]
是\Term{quasi-isomorphism}{拟同构}。于是断言 (a) 由 EGA $0_{\mathrm{III}}$ 13.2.3 的“\Term{topology}{拓扑}”变式得到。这就完成了 3.2.1 的证明。
\end{proof}

\textbf{3.2.2.} 若 $S$ 是\Term{topological-space}{拓扑空间}，以 $\operatorname{Parf}(S)$ 记 $(\mathcal O_S\otimes\mathbb C)$-\Term{module}{模}的\Term{perfect-complex}{完美复形}组成的\Term{triangulated-category}{三角范畴}，以 $K'(S)$ 记相应的“\Term{grothendieck-group}{Grothendieck 群}”\footnote{见第 I 讲 6.3。关于各 $K$ \Term{functor}{函子}的总体介绍，见第 IV 讲。}。
设 $S,X,E^\bullet$ 如 3.2.1；称元素
\[\operatorname{ind.an}(E^\bullet)=\operatorname{cl}(\mathrm Rp_*(E^\bullet))\in K'(S)\]
为 $E^\bullet$ 的\Term{analytic-index}{解析指标}。
\SourcePage{228}
若 $S$ 紧，$\operatorname{ind.an}(E^\bullet)$ \Term{canonical}{典范}地定义\Term{topology}{拓扑}学家通常使用的群 $K(S)$ 中的一个元素。事实上，\Term{generic}{一般}地，对每个\Term{topological-space}{拓扑空间} $S$，以 $\operatorname{Loclib}(S)$ 记\Term{locally-free-finite-type}{有限型局部自由} $(\mathcal O_S\otimes\mathbb C)$-\Term{module}{模}的\Term{additive-category}{加性范畴}，以 $K(S)$ 记\Term{triangulated-category}{三角范畴} $K^b(\operatorname{Loclib}(S))/K^{b,\phi}(\operatorname{Loclib}(S))$ 的“\Term{grothendieck-group}{Grothendieck 群}”（当 $S$ \Term{paracompact}{仿紧}时，该\Term{category}{范畴}等同于 $K^b(\operatorname{Loclib}(S))$，见第 II 讲 2.3.2(a)；因此在这种情况下，由第 I 讲 6.4，$K(S)$ \Term{canonical}{典范}地\Term{isomorphic}{同构}于\Term{additive-category}{加性范畴} $\operatorname{Loclib}(S)$ 的“\Term{grothendieck-group}{Grothendieck 群}”）。有自然\Term{homomorphism}{同态}
\[K(S)\longrightarrow K'(S).\]
但由第 II 讲 2.3.2，当 $S$ 紧时，它是\Term{isomorphism}{同构}，这就证明了我们的断言。更\Term{generic}{一般}地，$\operatorname{ind.an}(E^\bullet)$ 定义\Term{completion}{完备化}的群
\[\widehat K(S)=\varprojlim_A K(A)\]
中的一个元素，\Term{projective-limit}{投射极限}沿 $S$ 的紧\Term{subset}{子集} $A$ 取。

\textbf{3.2.3.} 编者不知道：如果不假定 $S$ \Term{locally-compact}{局部紧}，而只假定每个点有一个由\Term{paracompact}{仿紧}\Term{neighborhood}{邻域}组成的\Term{neighbourhood-basis}{基本邻域系}，\Term{finiteness-theorem}{有限性定理} 3.2.1 是否仍然成立。况且 $S$ \Term{locally-compact}{局部紧}的假设似乎并不很自然。事实上，固定阶数的\Term{elliptic-complex}{椭圆复形}的“泛\Term{local-scheme}{局部}参数”\Term{space}{空间}，是一个 Fréchet \Term{space}{空间}的一个（无限维）\Term{vector-subspace}{子空间}。将这一点说得更明确些。设 $V$ 是紧\Term{manifold}{流形}，$F=(F^i)_{i\in\mathbb Z}$ 是 $V$ 上\Term{finite-rank}{有限秩}复\Term{vector-bundle}{向量丛}的序列，除\Term{finite}{有限}项外均为零，$\underline n=(n_i)$ 是非负整数序列。以 $\operatorname{Comp.diff}^{\underline n}(F)$ 记“以 $F$ 为底”且阶数为 $\underline n$ 的\Term{differential-operator}{微分算子}\Term{complex}{复形}\Term{space}{空间}，即 $\prod_i\operatorname{Diff}^{n_i}(F^i,F^{i+1})$ 中由方程 $d^{i+1}d^i=0$ 定义的\Term{closed-subspace}{闭子空间}（\Term{space}{空间} $\prod_i\operatorname{Diff}^{n_i}(F^i,F^{i+1})$ 是 Fréchet \Term{space}{空间}）。
\SourcePage{229}
其中由阶数为 $\underline n$ 的\Term{elliptic-complex}{椭圆复形}组成的\Term{subset}{子集} $\operatorname{Ell}^{\underline n}(F)$ 是\Term{open-subset}{开集}；正是这个\Term{space}{空间}（\Term{generic}{一般}并不\Term{locally-compact}{局部紧}）起着泛\Term{local-scheme}{局部}\Term{parameter-space}{参数空间}的\Term{action}{作用}。

\textbf{3.2.4.} 不过，即使上面提到的 3.2.1 的推广不可行，只假定\Term{fibration}{纤维化} $p$ \Term{proper}{固有}且\Term{locally-trivial}{局部平凡}，仍可给 $E^\bullet$ 配上 $K'(S)$ 中的一个元素，称它为 $E^\bullet$ 的\Term{analytic-index}{解析指标}是恰当的。事实上，为简便起见\footnote{在\Term{generic}{一般}情形中，使用 [12] 的“\Term{weak-hilbert-bundle}{弱 Hilbert 丛}”，即 Hilbert \Term{space}{空间}的\Term{continuous-field}{连续场}的概念。}，假定处于“\Term{local-scheme}{局部}”情形：$X=S\times V$（$V$ 为紧\Term{manifold}{流形}），$E^\bullet$ 由一个从 $S$ 到 $\operatorname{Ell}^{\underline n}(F)$ 的\Term{continuous-map}{连续映射}给出（3.2.3）。3.2.1 证明中的断言 (a')、(b) 实际上不需要对 $S$ 作任何假设；所以采用该处记号，置
\[\operatorname{ind.an}(E^\bullet)=\operatorname{cl}(\mathcal F_r^\bullet)\in K'(S)\]
是合理的，该元素与 $r$ 无关。

\textbf{3.2.5.} 假设 $S$ \Term{paracompact}{仿紧}，以 $H$ 记一个 Hilbert \Term{dimension}{维数}无限的复 Hilbert \Term{space}{空间}，以 $I(H)$ 记 $L(H,H)$ 中由“有指标”的\Term{operator}{算子}组成的\Term{vector-subspace}{子空间}。那么各\Term{complex}{复形} $F_r^\bullet$ 定义了群 $[S,I(H)]$ 中的一个\Term{canonical}{典范}元素（[12]）；这个群由从 $S$ 到 $I(H)$ 的\Term{continuous-map}{连续映射}的\Term{homotopy-class}{同伦类}组成。但当 $S$ 不紧时（甚至 $S$ \Term{locally-compact}{局部紧}时），编者不知道群 $[S,I(S)]$ 与 $K'(S)$ 之间是什么关系（是否存在从其中一个到另一个的自然\Term{arrow}{箭头}？它们是否\Term{isomorphic}{同构}？），因而更不知道如此定义的两个指标之间的关系。

关于在 3.2.1 假设下对 Atiyah--Singer 指标\Term{theorem}{定理}的推广，参见 Shih 最近的工作 [17]，或 M.F. Atiyah 与 I.M. Singer 正在准备的一篇论文。

\SourcePage{230}
\SGAsection{参考文献}{第 II 讲参考文献}
\begin{enumerate}
\item[{[1]}] Artin, M.、Grothendieck, A. 与 Verdier, J.-L.，\emph{\Term{scheme}{概形}的\Term{etale-cohomology}{平展上同调}}（Cohomologie étale des schémas），将由 North Holl. Pub. Comp. 出版（引作 SGA 4）。
\item[{[2]}] Atiyah, M. 与 Segal, G.，未公开的文稿（papiers secrets）。
\item[{[3]}] Cartan, H.，\emph{复函数与复\Term{manifold}{流形}；\Term{coherent}{凝聚}解析层}（Funzioni e varietà complesse, Faisceaux analytiques cohérents），Centro Internazionale Matematico Estivo 课程，1964，Roma。
\item[{[4]}] Cartan, H. 与 Schwartz, L.，\emph{Atiyah--Singer \Term{theorem}{定理}}（Théorème d'Atiyah--Singer），1963--64 年讨论班，Institut Henri Poincaré，Paris。
\item[{[5]}] Frisch, J.，\emph{半解析集上的解析函数}（Fonctions analytiques sur un ensemble semi-analytique），CRAS Paris 短文，第 260 卷，第 2974--2976 页（1965）。
\item[{[6]}] Godement, R.，\emph{层论}（Théorie des Faisceaux），Hermann（1958）（引作 T.F.）。
\item[{[7]}] Grothendieck, A.，\emph{\Term{coherent-sheaf}{凝聚层}的\Term{local-cohomology}{局部上同调}及\Term{local-scheme}{局部}和整体 Lefschetz \Term{theorem}{定理}}（Cohomologie locale des faisceaux cohérents et théorèmes de Lefschetz locaux et globaux），IHES 1962 年\Term{algebraic-geometry}{代数几何}讨论班，将由 North Holl. Pub. Comp. 出版（引作 SGA 2）。
\item[{[8]}] Grothendieck, A.，\emph{\Term{analytic-geometry}{解析几何}基础}（Fondements de la Géométrie Analytique），Cartan 讨论班，1959--60，Institut Henri Poincaré，Paris。
\item[{[9]}] Hartshorne, R.，\emph{\Term{residue}{留数}与\Term{duality}{对偶}}（Residues and Duality），Lecture Notes，Springer，第 20 号，引作 [H]。
\item[{[10]}] Houzel, C.，\emph{\Term{local-scheme}{局部}\Term{analytic-geometry}{解析几何}}（Géométrie Analytique Locale），Cartan 讨论班，1959--60，Institut Henri Poincaré，Paris。
\item[{[11]}] Illusie, L.，\emph{丛\Term{complex}{复形}与\Term{functor}{函子} $K$ 的定义}（Complexes de fibrés et définitions du foncteur $K$），Shih Weishu 讨论班，IHES，1964--65。
\item[{[12]}] Illusie, L.，\emph{Banach 丛的\Term{direct-quasi-acyclic-complex}{直接拟无环复形}}（Complexes quasi-acycliques directs de fibrés banachiques），CRAS Paris 短文，第 260 卷，第 6499--6502 页。
\item[{[13]}] Kiehl, R.，\emph{非阿基米德函数论中的\Term{theorem}{定理} A 与\Term{theorem}{定理} B}（Theorem A und Theorem B in der nichtarchimedischen Funktionentheorie），Invent. math. 2，第 256--273 页（1967）。
\SourcePage{231}
\item[{[14]}] Lang, S.，\emph{\Term{differentiable-manifold}{微分流形}导论}（Introduction to differentiable manifolds）。
\item[{[15]}] Łojasiewicz, S.，Ann. della Scuola Norm. Sup. Pisa，第三辑，18，第 IV 分册，1964。
\item[{[16]}] Palais, R.，\emph{关于指标\Term{theorem}{定理}的 Princeton 讨论班报告（1963--64）}（Exposés au Séminaire Princeton 1963--64 sur le théorème de l'index）。
\item[{[17]}] Shih, W.，\emph{\Term{fiber}{纤维}配边与\Term{elliptic-operator}{椭圆微分算子}族的指标}（Fiber cobordism and the index of a family of elliptic differential operators）。
\item[{[18]}] Verdier, J.-L.，\emph{\Term{abelian-category}{阿贝尔范畴}的\Term{derived-category}{导出范畴}}（Les catégories dérivées des catégories abéliennes），将由 North Holl. Pub. Comp. 出版。
\end{enumerate}
